% GENERATED FILE. Edit canonical lessons, then run npm run generate:books.
% canonical-source-hash: fnv1a64:59b76ec4adf205cb
% canonical-lessons: ZH-C116-qianbao, ZH-C116-huzhao, ZH-C116-shoubiao, ZH-C116-yusan, ZH-C116-diu, ZH-R116-review-111-116

\chapter{Things I Carry}
\label{ch:zh-things-i-carry}
\hlchaptermodality{116}

\begin{chapteropening}
\textbf{By the end of this chapter:} \emph{I can say I have my wallet, passport, watch and umbrella, and say I have lost something.}

Complete the last lesson of chapter 116: I can say I have my wallet, passport, watch and umbrella, and say I have lost something.
\end{chapteropening}

\section[qiánbāo]{qiánbāo — wallet; purse}

\label{lesson:ZH-C116-qianbao}



\begin{quote}
Before the new one: say the Mandarin for to pay attention, then the Mandarin for dangerous.
\end{quote}

\subsection*{You'll want to know: qiánbāo}
\emph{qiánbāo} — \textquotedblleft{}wallet; purse\textquotedblright{}.

\begin{quote}
\emph{wǒ de qiánbāo zài bāo lǐ} — My wallet's in my bag.
\end{quote}

Its characters come later; for now it is a word for the ear.

\begin{cousinweb}
\emph{qiánbāo} is built from qián 'money' and bāo 'to wrap', both inherited from Middle Chinese, \emph{dzjen} and \emph{paew} (Guangyun, in Baxter's transcription).
\end{cousinweb}

\begin{grammarlens}[title={What you've built}]
One more word for what I have and can do.
\end{grammarlens}

\subsection*{Your turn}
\begin{itemize}
\raggedright
  \item \emph{Say it:} \emph{qiánbāo}
  \item \emph{Say it:} \emph{qiánbāo}, once more
  \item \emph{Say it:} \emph{wǒ de qiánbāo zài bāo lǐ} — My wallet's in my bag.
  \item \emph{Recall:} say the Mandarin for to pay attention, then the Mandarin for dangerous, then say \emph{qiánbāo} again
\end{itemize}

\begin{tcolorbox}[breakable,colback=teal!4,colframe=teal!35!black,title={Before you move on}]
What does qiánbāo mean? (\textquotedblleft{}Wallet; purse\textquotedblright{}.) Say it once more.
\end{tcolorbox}
\section[hùzhào]{hùzhào — passport}

\label{lesson:ZH-C116-huzhao}



\begin{quote}
Before the new one: say the Mandarin for dangerous, then the Mandarin for wallet.
\end{quote}

\subsection*{You'll want to know: hùzhào}
\emph{hùzhào} — \textquotedblleft{}passport\textquotedblright{}.

\begin{quote}
\emph{nǐ dài hùzhào le ma} — Have you brought your passport?
\end{quote}

Its characters come later; for now it is a word for the ear.

\begin{cousinweb}
\emph{hùzhào} is built from hù 'to protect' and zhào 'bright', both inherited from Middle Chinese, \emph{huH} and \emph{tsyewH} (Guangyun, in Baxter's transcription). The passport is a modern sense of the pair, and how it came to it is not settled.
\end{cousinweb}

\begin{grammarlens}[title={What you've built}]
One more word for what I have and can do.
\end{grammarlens}

\subsection*{Your turn}
\begin{itemize}
\raggedright
  \item \emph{Say it:} \emph{hùzhào}
  \item \emph{Say it:} \emph{hùzhào}, once more
  \item \emph{Say it:} \emph{nǐ dài hùzhào le ma} — Have you brought your passport?
  \item \emph{Recall:} say the Mandarin for dangerous, then the Mandarin for wallet, then say \emph{hùzhào} again
\end{itemize}

\begin{tcolorbox}[breakable,colback=teal!4,colframe=teal!35!black,title={Before you move on}]
What does hùzhào mean? (\textquotedblleft{}Passport\textquotedblright{}.) Say it once more.
\end{tcolorbox}
\section[shǒubiǎo]{shǒubiǎo — wristwatch}

\label{lesson:ZH-C116-shoubiao}



\begin{quote}
Before the new one: say the Mandarin for wallet, then the Mandarin for passport.
\end{quote}

\subsection*{You'll want to know: shǒubiǎo}
\emph{shǒubiǎo} — \textquotedblleft{}wristwatch\textquotedblright{}.

\begin{quote}
\emph{nǐ de shǒubiǎo hěn hǎokàn} — Your watch is lovely.
\end{quote}

Its characters come later; for now it is a word for the ear.

Two third tones: shǒu rises before biǎo, and is still written shǒu.

\begin{cousinweb}
\emph{shǒubiǎo} is built from \zh{手} 'hand' and biǎo, both inherited from Middle Chinese, \emph{syuwX} and \emph{pjewX} (Guangyun, in Baxter's transcription); the dictionary glosses the second 'to make clear' and 'an outer garment'. The watch is a modern sense, and how biǎo came to it is not settled.
\end{cousinweb}

\begin{grammarlens}[title={What you've built}]
One more word for what I have and can do.
\end{grammarlens}

\subsection*{Your turn}
\begin{itemize}
\raggedright
  \item \emph{Say it:} \emph{shǒubiǎo}
  \item \emph{Say it:} \emph{shǒubiǎo}, once more
  \item \emph{Say it:} \emph{nǐ de shǒubiǎo hěn hǎokàn} — Your watch is lovely.
  \item \emph{Recall:} say the Mandarin for wallet, then the Mandarin for passport, then say \emph{shǒubiǎo} again
\end{itemize}

\begin{tcolorbox}[breakable,colback=teal!4,colframe=teal!35!black,title={Before you move on}]
What does shǒubiǎo mean? (\textquotedblleft{}Wristwatch\textquotedblright{}.) Say it once more.
\end{tcolorbox}
\section[yǔsǎn]{yǔsǎn — umbrella}

\label{lesson:ZH-C116-yusan}



\begin{quote}
Before the new one: say the Mandarin for passport, then the Mandarin for wristwatch.
\end{quote}

\subsection*{You'll want to know: yǔsǎn}
\emph{yǔsǎn} — \textquotedblleft{}umbrella\textquotedblright{}.

\begin{quote}
\emph{xiàyǔ le nǐ dài yǔsǎn le ma} — It's raining; have you got an umbrella?
\end{quote}

Its characters come later; for now it is a word for the ear.

Two third tones: yǔ rises before sǎn, and is still written yǔ.

\begin{cousinweb}
\emph{yǔsǎn} is built from \zh{雨} 'rain' and sǎn, both inherited from Middle Chinese, \emph{hjuX} and \emph{sanX} (Guangyun, in Baxter's transcription); the dictionary glosses the second as a canopy.
\end{cousinweb}

\begin{grammarlens}[title={What you've built}]
One more word for what I have and can do.
\end{grammarlens}

\subsection*{Your turn}
\begin{itemize}
\raggedright
  \item \emph{Say it:} \emph{yǔsǎn}
  \item \emph{Say it:} \emph{yǔsǎn}, once more
  \item \emph{Say it:} \emph{xiàyǔ le nǐ dài yǔsǎn le ma} — It's raining; have you got an umbrella?
  \item \emph{Recall:} say the Mandarin for passport, then the Mandarin for wristwatch, then say \emph{yǔsǎn} again
\end{itemize}

\begin{tcolorbox}[breakable,colback=teal!4,colframe=teal!35!black,title={Before you move on}]
What does yǔsǎn mean? (\textquotedblleft{}Umbrella\textquotedblright{}.) Say it once more.
\end{tcolorbox}
\section[diū]{diū — to lose (a thing)}

\label{lesson:ZH-C116-diu}



\begin{quote}
Before the new one: say the Mandarin for wristwatch, then the Mandarin for umbrella.
\end{quote}

\subsection*{You'll want to know: diū}
\emph{diū} — \textquotedblleft{}to lose (a thing)\textquotedblright{}.

\begin{quote}
\emph{wǒ de yàoshi diū le} — I've lost my keys.
\end{quote}

Its characters come later; for now it is a word for the ear.

\begin{cousinweb}
\emph{diū} is the everyday verb for losing a thing. No dictionary we use records an older form, so its history is not settled.
\end{cousinweb}

\begin{grammarlens}[title={What you've built}]
One more word for what I have and can do.
\end{grammarlens}

\subsection*{Your turn}
\begin{itemize}
\raggedright
  \item \emph{Say it:} \emph{diū}
  \item \emph{Say it:} \emph{diū}, once more
  \item \emph{Say it:} \emph{wǒ de yàoshi diū le} — I've lost my keys.
  \item \emph{Recall:} say the Mandarin for wristwatch, then the Mandarin for umbrella, then say \emph{diū} again
\end{itemize}

\begin{tcolorbox}[breakable,colback=teal!4,colframe=teal!35!black,title={Before you move on}]
What does diū mean? (\textquotedblleft{}To lose (a thing)\textquotedblright{}.) Say it once more.
\end{tcolorbox}
\section[(dialogue)]{Looking back: from In Front and Behind to Things I Carry}

\label{lesson:ZH-R116-review-111-116}



\begin{quote}
Say the Mandarin for in front and at the back.
\end{quote}

\subsection*{Your turn}
\begin{itemize}
\raggedright
  \item \emph{Say it:} in front, at the back, nearby, crossroads, to turn a corner
  \item \emph{Say it:} exit, entrance, to push, to pull, stairs
  \item \emph{Say it:} to slice, to boil, to stir-fry, to roast, to fry
  \item \emph{Say it:} to pay the bill, to pay, cash, to pay by card, to give change
  \item \emph{Say it:} to be careful, to smoke, to forbid, to pay attention, dangerous
  \item \emph{Say it:} wallet, passport, wristwatch, umbrella, to lose
  \item \emph{Read it:} \zh{边} \zh{出} \zh{心} — each one aloud, then a word that holds it
\end{itemize}

\begin{tcolorbox}[breakable,colback=teal!4,colframe=teal!35!black,title={Before you move on}]
In front? (\textbf{\zh{前边}}, \emph{qiánbian}.) At the back? (\emph{hòubian}.) Nearby? (\emph{fùjìn}.) Crossroads? (\emph{lùkǒu}.) To turn a corner? (\emph{guǎi}.) Exit? (\textbf{\zh{出口}}, \emph{chūkǒu}.) Entrance? (\emph{rùkǒu}.) To push? (\emph{tuī}.) To pull? (\emph{lā}.) Stairs? (\emph{lóutī}.) To slice? (\emph{qiē}.) To boil? (\emph{zhǔ}.) To stir-fry? (\emph{chǎo}.) To roast? (\emph{kǎo}.) To fry? (\emph{jiān}.) To pay the bill? (\emph{mǎidān}.) To pay? (\emph{fù}.) Cash? (\emph{xiànjīn}.) To pay by card? (\emph{shuākǎ}.) To give change? (\emph{zhǎoqián}.) To be careful? (\textbf{\zh{小心}}, \emph{xiǎoxīn}.) To smoke? (\emph{xīyān}.) To forbid? (\emph{jìnzhǐ}.) To pay attention? (\emph{zhùyì}.) Dangerous? (\emph{wēixiǎn}.) Wallet? (\emph{qiánbāo}.) Passport? (\emph{hùzhào}.) Wristwatch? (\emph{shǒubiǎo}.) Umbrella? (\emph{yǔsǎn}.) To lose? (\emph{diū}.)
\end{tcolorbox}
